\documentclass[11pt]{article}
\usepackage{amsmath, amssymb, amsthm, enumerate}
\usepackage[margin=1in, top=1.2in, bottom=1.2in]{geometry}
\usepackage{parskip}
\usepackage[colorlinks=true, linkcolor=blue, citecolor=blue, urlcolor=blue]{hyperref}
\usepackage{cleveref}
\theoremstyle{definition}
\newtheorem{definition}{Definition}[section]
\newtheorem{theorem}{Theorem}[section]
\newtheorem{lemma}{Lemma}[section]
\newtheorem{proposition}{Proposition}[section]
\newtheorem{corollary}{Corollary}[section]
\theoremstyle{remark}
\newtheorem{remark}{Remark}[section]
\newtheorem{example}{Example}[section]
\setlength{\parindent}{0pt}
\setlength{\parskip}{6pt}
\begin{document}

\begin{center}
{\Large\textbf{homework_1}}\\
\vspace{0.3cm}
{\normalsize\today}
\end{center}
\vspace{0.5cm}


% --- Page 1 ---
\begin{enumerate}
    \item If $R$ is a ring, show that $M_n(R)$ is a ring under the usual operations of matrix addition and multiplication.

    \item Let $(G, +)$ be an Abelian group. An \textbf{endomorphism} of $G$ is a group homomorphism $\varphi : G \to G$. Let $\text{End}(G)$ denote the set of all endomorphisms on $G$. Given $\varphi, \psi \in \text{End}(G)$, define
    \[
    (\varphi + \psi)(g) := \varphi(g) + \psi(g), \quad \text{and} \quad (\varphi \times \psi)(g) := \varphi(\psi(g)).
    \]
    \begin{enumerate}
        \item[(a)] Show that $+$ and $\times$ are well-defined binary operations on $\text{End}(G)$.
        \item[(b)] Show that $(\text{End}(G), +, \times)$ is a ring.
    \end{enumerate}

    \item Let $R$ be a ring and $a, b \in R$. Show that
    \begin{enumerate}
        \item[(a)] $0 \cdot a = a \cdot 0 = 0$.
        \item[(b)] For any $n \in \mathbb{Z}$, $(na)b = n(ab) = a(nb)$. [Hint: Use Induction.]
        \item[(c)] If $R$ is unital, then the multiplicative identity is unique. Moreover, $(-1) \cdot a = (-a)$ for all $a \in R$.
    \end{enumerate}

    \item Let $R$ be a commutative ring and let $S$ and $T$ be two subsets of $R$. Define
    \[
    ST := \bigcup_{n=1}^{\infty} \left\{ \sum_{i=1}^n s_i t_i : s_i \in S, t_i \in T \right\}.
    \]
    (In other words, elements in $ST$ are finite sums of the form $\sum_{i=1}^n s_i t_i$, where $s_i \in S$ and $t_i \in T$, while the number $n$ may vary.)
    \begin{enumerate}
        \item[(a)] If $S$ and $T$ are both subrings of $R$, prove that $ST$ is a subring.
        \item[(b)] If $S$ and $T$ are both ideals, then prove that $ST$ is an ideal.
    \end{enumerate}

    \item Fix $m, n \in \mathbb{Z}$ and let $d := \gcd(m, n)$. If $I$ denoted the ideal generated by $\{m, n\}$, then show that $I = (d)$, the principal ideal generated by $d$ (See Definition 5.3).
\end{enumerate}

% --- Page 2 ---
Extra Problems

You need not submit these problems; they are here only for practice.

**Definition:** A function $f : \mathbb{R} \to \mathbb{R}$ is said to have **compact support** if there exists $a, b \in \mathbb{R}$ such that $f(x) = 0$ for all $x \notin [a, b]$.

Define
\[ C_c(\mathbb{R}) := \{ f : \mathbb{R} \to \mathbb{C} : f \text{ is continuous and has compact support} \}. \]

Equip $C_c(\mathbb{R})$ with the pointwise operations as in Part (v) of Example 1.2 in the notes. You need not check that $C_c(\mathbb{R})$ is a ring (this is easy to do).

\begin{enumerate}
    \item Prove that $C_c(\mathbb{R})$ is not unital.

    \item Let $R$ be a ring, $S$ be a subring of $R$ and let $I$ be an ideal of $R$.
    \begin{enumerate}
        \item Prove that $S + I := \{ s + a : s \in S, a \in I \}$ is a subring of $R$.
        \item Prove that $I \trianglelefteq S + I$.
        \item Prove that $S \cap I \trianglelefteq S$.
        \item Prove that
        \[
        \frac{S + I}{I} \cong \frac{S}{S \cap I}
        \]
        \textit{[Hint: Define $\varphi : S \to (S + I)/I$ by $s \mapsto s + I$.]}
    \end{enumerate}

    \item Let $D \in \mathbb{Q}$ be a rational number that is not a perfect square. Define
    \[
    R := \{ a + b\sqrt{D} : a, b \in \mathbb{Q} \}.
    \]
    \begin{enumerate}
        \item Show that $R$ is a subring of $\mathbb{C}$.
        \item If $(a + b\sqrt{D}) \neq 0$, then show that $a^2 - b^2 D \neq 0$.
        \item Use Part (b) to prove that $R$ is a field.
        This field is denoted $\mathbb{Q}(\sqrt{D})$, and is called a \textbf{quadratic field}.
    \end{enumerate}

    \item For a unital ring $R$, define the \textbf{center} of $R$ to be
    \[
    \mathcal{Z}(R) := \{ a \in R : ab = ba \text{ for all } b \in R \}.
    \]
    Show that $\mathcal{Z}(R)$ is a subring of $R$.
\end{enumerate}

% --- Page 3 ---
Given the content in the image, 
---

\begin{enumerate}
    \item Given: $R$ is a ring.

    Show: $M_n(\mathbb{C} R)$ is a ring under $+$, matrix multiplication.

    \begin{enumerate}
        \item[(iii)] Fix $n \in \mathbb{N}$, show: $M_n(\mathbb{C} R)$ is a ring under addition and multiplication of real numbers is $C = AB$ s.t.
        \[
        c_{ij} = \sum_{k=1}^n a_{ik} b_{kj} \quad (\text{Check})
        \]

        \textbf{Associativity of multiplication i.e.}
        \[
        A \cdot (B \cdot C) = A \cdot D \quad \text{where} \quad d_{ij} = \sum_{k=1}^n b_{ik} c_{kj}
        \]
        \[
        \Rightarrow A \cdot (B \cdot C) = \left( \sum_{k=1}^n a_{ik} d_{kj} \right)_{n \times n}
        \]
        \[
        = \left( \sum_{k=1}^n a_{ik} \left( \sum_{l=1}^n b_{il} c_{lj} \right) \right)_{n \times n}
        \]
        \[
        = \left( \sum_{l=1}^n \left( \sum_{k=1}^n a_{ik} b_{kl} \right) c_{lj} \right)_{n \times n}
        \]
        \[
        = (A \cdot B) \cdot C
        \]

        \textbf{Distributive Property}
        \[
        AB + AC = \left( \sum_{k=1}^n a_{ik} b_{kj} + \sum_{k=1}^n a_{ik} c_{kj} \right)_{n \times n}
        \]
        \[
        = \left( \sum_{k=1}^n a_{ik} (b_{kj} + c_{kj}) \right)_{n \times n}
        \]
        \[
        = A (B + C)
        \]

        (We can similarly show $(A + B) C = AC + BC$.)

        \begin{proof}
        $\square$
        \end{proof}

        \item[(iv)] Conclude that if $R$ is a ring, then $M_n(R)$ is a ring.
    \end{enumerate}
\end{enumerate}

% --- Page 4 ---
Given the content in the image, 
---

Given $(G, +)$ is an Abelian group,

\[
\text{End}(G) = \left\{ \varphi : G \to G \mid \varphi \text{ is a group homomorphism} \right\}
\]

For $\varphi, \psi \in \text{End}(G)$,

\[
(\varphi + \psi)(g) = \varphi(g) + \psi(g)
\]

\[
(\varphi \times \psi)(g) = \varphi(\psi(g))
\]

(a)

\begin{proof}
We show the operations given above are well-defined.

\textbf{Well-definedness of $+$:}
We need to show $\varphi + \psi \in \text{End}(G)$. Consider $g, h \in G$,

\[
(\varphi + \psi)(g + h) = \varphi(g + h) + \psi(g + h) = \varphi(g) + \varphi(h) + \psi(g) + \psi(h)
\]

Using the Abelian property,

\[
= \varphi(g) + \psi(g) + \varphi(h) + \psi(h) = (\varphi + \psi)(g) + (\varphi + \psi)(h)
\]

Thus, $\varphi + \psi \in \text{End}(G)$.

\textbf{Well-definedness of $\times$:}
We need to show $\varphi \times \psi \in \text{End}(G)$. Consider $g, h \in G$,

\[
(\varphi \times \psi)(g + h) = \varphi(\psi(g + h)) = \varphi(\psi(g) + \psi(h)) = \varphi(\psi(g)) + \varphi(\psi(h))
\]

\[
= (\varphi \times \psi)(g) + (\varphi \times \psi)(h)
\]

Thus, $\varphi \times \psi \in \text{End}(G)$.
\end{proof}

(b)

\begin{theorem}\label{thm:end_group_ring}
Prove $(\text{End}(G), +, \times)$ is a ring.
\end{theorem}

\begin{proof}
Since $G$ is an Abelian group, $(\text{End}(G), +)$ is an Abelian group.

\textbf{Associativity of $\times$:}
The operation $\times$ is the composition of functions, which is associative.

\textbf{Distributive Property:}
For $\varphi, \psi_1, \psi_2 \in \text{End}(G)$ and $g \in G$,

\[
\varphi \times (\psi_1 + \psi_2)(g) = \varphi((\psi_1 + \psi_2)(g)) = \varphi(\psi_1(g) + \psi_2(g))
\]

\[
= \varphi(\psi_1(g)) + \varphi(\psi_2(g)) = (\varphi \times \psi_1)(g) + (\varphi \times \psi_2)(g)
\]

Thus, $(\text{End}(G), +, \times)$ is a ring.
\end{proof}

% --- Page 5 ---
---
\[
(\varphi_1 + \varphi_2) \times \varphi(g) = (\varphi_1 + \varphi_2)(\varphi(g)) = \varphi_1(\varphi(g)) + \varphi_2(\varphi(g)) = \varphi_1 \times \varphi(g) + \varphi_2 \times \varphi(g)
\]

\begin{enumerate}
    \item Given: $(R, +, \cdot)$ is a ring; $a, b \in R$

    \begin{enumerate}
        \item[a)]
        \begin{proof}
            Show: $O \cdot a = a \cdot O = O$

            \[
            O \cdot a = (a - a) \cdot a = a^2 - a^2 = O
            \]
            \[
            a \cdot O = a \cdot (a - a) = a^2 - a^2 = O
            \]
        \end{proof}

        \item[b)]
        \begin{theorem}\label{thm:scalar_multiplication}
            Show: $(na)b = n(ab) = a(nb)$
        \end{theorem}
        \begin{proof}
            Base case $(n=1)$: Trivial.

            Induction Hypothesis: $(na)b = n(ab) = a(nb)$

            \[
            (a + \underbrace{a + \cdots + a}_{n+1 \text{ times}})b = (na + a)b = nab + ab = (n+1)ab
            \]
            \[
            a(\underbrace{n \text{ times } b}_{n+1 \text{ times}}) = a(nb + b) = nab + ab
            \]
        \end{proof}

        \item[c)]
        \begin{theorem}\label{thm:unique_multiplicative_identity}
            If $R$ is unital, prove: $1_R$ is multiplicative identity is unique.
        \end{theorem}
        \begin{proof}
            Let $1_A$ and $1_B$ be multiplicative identities in $R$. Then,
            \[
            1_A = 1_A \cdot 1_B = 1_B
            \]
            Thus, the multiplicative identity is unique.

            Show: $(-1) \cdot a = -a$ for any $a \in R$.

            Proof:
            \[
            a + (-1) \cdot a = 1 \cdot a + (-1) \cdot a = (1 + (-1)) \cdot a = 0
            \]
            Therefore, $(-1) \cdot a$ is the additive inverse of $a$, so $(-1) \cdot a = -a$.
        \end{proof}
    \end{enumerate}

    \item Given: $R$ is a commutative ring,
    \[
    S, T \subseteq R
    \]
    \[
    ST = \bigcup_{n=1}^{\infty} \left\{ \sum_{i=1}^n s_i t_i \mid s_i \in S, t_i \in T \right\}
    \]

    \begin{enumerate}
        \item[a)]
        \begin{theorem}\label{thm:subring_product}
            If $S$ and $T$ are both subrings, show: $ST$ is a subring.
        \end{theorem}
        \begin{proof}
            Assume $S$ and $T$ to be non-empty.
        \end{proof}
    \end{enumerate}
\end{enumerate}
---

% --- Page 6 ---
---

Let $a, b \in ST$ such that $a = \sum_{i=1}^{n_1} s_i t_i$ and $b = \sum_{i=1}^{n_2} s_i' t_i'$.

Without loss of generality, assume $n_2 > n_1$. Then,
\[
b - a = \sum_{i=1}^{n_2} s_i' t_i' - \sum_{i=1}^{n_1} s_i t_i = \sum_{i=n_1+1}^{n_2} s_i' t_i' + \sum_{i=1}^{n_1} (s_i' t_i' - s_i t_i).
\]

By the Abelian property, we can rearrange the terms:
\[
\sum_{i=n_1+1}^{n_2} s_i' t_i' + \sum_{i=1}^{n_1} (s_i' t_i' - s_i t_i) = \sum_{i=1}^{n_2} s_i'' t_i''
\]
where $i = 1$ to $i = n_1$ is set to $0$ for the first entry.

Thus, $(ST, +) \prec (R, +)$.

\begin{proof}
Check the ring property:
\[
\left(\sum_{i=1}^{n_1} s_i t_i\right) \left(\sum_{j=1}^{n_2} s_j' t_j'\right).
\]

Without loss of generality, assume $n_1 < n_2$. Now set all terms to zero in the first entry after $i > n_1$:
\[
\left(\sum_{i=1}^{n_2} s_i t_i\right) \left(\sum_{j=1}^{n_2} s_j' t_j'\right).
\]

Every term will be of the form $s_i s_j' t_i t_j'$ for some $1 \leq i, j \leq n_2$. As $R$ is a commutative ring, we get $s_i s_j' = s_j' s_i$ and $t_i t_j' = t_j' t_i$. Since $S, T$ are subrings, $s_i s_j' \in S$ and $t_i t_j' \in T$ respectively.
\end{proof}

\begin{definition}
If $S, T$ are both ideals, then $ST$ is an ideal.
\end{definition}

\begin{proof}
Consider $\sum_{i=1}^{n} s_i t_i \in ST$ and $r \in R$. Then,
\[
r \left(\sum_{i=1}^{n} s_i t_i\right) = \sum_{i=1}^{n} (r s_i) t_i.
\]
As $S$ is an ideal, we get $r s_i \in S$ for all $i \leq n$. Thus,
\[
\sum_{i=1}^{n} (r s_i) t_i \in ST.
\]

Similarly,
\[
\left(\sum_{i=1}^{n} s_i t_i\right) r = \sum_{i=1}^{n} s_i (t_i r).
\]
As $T$ is an ideal, we get $t_i r \in T$ for all $i \leq n$. Thus,
\[
\sum_{i=1}^{n} s_i (t_i r) \in ST.
\]
\end{proof}

---

% --- Page 7 ---
---

Given $m, n \in \mathbb{Z}$ and $d = \gcd(m, n)$,

\begin{definition}\label{def:ideal_generated}
Let $I$ be the ideal generated by $\{m, n\}$.
\end{definition}

\begin{theorem}\label{thm:ideal_gcd}
Show that $I = (d)$.
\end{theorem}

\begin{proof}
Consider $C(X) = \bigcap_{J \in \mathcal{F}} J$ where
\[
\mathcal{F} = \{ J \triangleleft \mathbb{Z} \mid X \subseteq J \}.
\]
If $a \in C(X)$, then $a \in J$ for every $J \in \mathcal{F}$.

Since $J$ is an ideal of $\mathbb{Z}$, it must also be a subgroup of $\mathbb{Z}$. Thus, $\exists n_J \in \mathbb{Z}$ such that $J = n_J \mathbb{Z}$.

If $X \subseteq n_J \mathbb{Z}$, then $n_J \mid x$ and $n_J \mid y$ for $x, y \in X$.

Since $d = \gcd(x, y)$, we have $d \mid n_J$.

Thus, $n_J \mathbb{Z} \supseteq d \mathbb{Z}$ for every $J \in \mathcal{F}$.

Therefore, $C(X) \supseteq d \mathbb{Z}$.

\begin{claim}
$d \mathbb{Z}$ is the smallest ideal containing $X$.
\end{claim}

Consider $n \mathbb{Z} \triangleleft \mathbb{Z}$ such that $X \subseteq n \mathbb{Z}$ and $n \mathbb{Z} \subseteq d \mathbb{Z}$.

Let $nk_1 = x$ and $nk_2 = y$ for some $k_1, k_2 \in \mathbb{Z}$.

By Bézout's identity, $nk_1 r_1 + nk_2 r_2 = d$ for some $r_1, r_2 \in \mathbb{Z}$.

Thus, $n \mid d$.

Hence, $d \mathbb{Z} \subseteq n \mathbb{Z}$ and $n \mathbb{Z} = d \mathbb{Z}$.

\end{proof}

---

% --- Page 8 ---
---

\begin{definition}\label{def:cc}
Given the set
\[
C_c(\mathbb{R}) := \{ f : \mathbb{R} \to \mathbb{C} \mid f \text{ is continuous and has compact support} \},
\]
where compact support means $\exists a, b \in \mathbb{R}$ such that $f(x) = 0$ whenever $x \notin [a, b]$,
and $(+, \times)$ are the same as usual addition and multiplication.
\end{definition}

\begin{remark}\label{rem:ccring}
Observe that $C_c(\mathbb{R})$ is a ring, and it is non-empty.
\end{remark}

\begin{theorem}\label{thm:nonunit}
Show that $C_c(\mathbb{R})$ is \textbf{not} unital.
\end{theorem}

\begin{proof}
Assume there is $g \in C_c(\mathbb{R})$ such that $f \times g = f$ for any $f \in C_c(\mathbb{R})$.

If $g(x_0) \neq 0$ and $x_0 \in [a, b]$, choose $f \in C_c(\mathbb{R})$ such that $f(x_0) \neq 0$ and $x_0 \in [a', b']$ where $[a, b] \subset [a', b']$. Then observe at $x_0 = \frac{a' + a}{2}$,
\[
g(x_0) = 0 \quad \text{and} \quad f(x_0) \neq 0 \implies f \times g(x_0) \neq 0.
\]
This is a contradiction. Hence, no unit exists in $C_c(\mathbb{R})$.
\end{proof}

% --- Page 9 ---
Given the content in the image, 
---

\begin{enumerate}
    \item Given: $R$ is a ring, $S$ be a subring of $R$ and $I \trianglelefteq R$.

    \begin{enumerate}
        \item[(a)]
        \begin{proof}
            Prove: $S + I := \{st + a \mid s \in S, a \in I\}$ is a subring.

            \textbf{Proof.} Observe, $S + I$ is non-empty, then consider $a', b' \in S + I$ where $a' = s_1 + a$ and $b' = s_2 + b$. Now,
            \[
            a' - b' = s_1 + a - (s_2 + b) = (s_1 - s_2) + (a - b) \in S + I
            \]
            where $s_1 - s_2 \in S$ and $a - b \in I$. This implies $(S + I, +) \leq (R, +)$.

            Now, let $a', b' \in S + I$ where $a' = s_1 + a$, $b' = s_2 + b$. Then,
            \[
            a' \cdot b' = (s_1 + a)(s_2 + b) = s_1s_2 + as_2 + s_1b + ab
            \]
            where $s_1s_2 \in S$ and $as_2 + s_1b + ab \in I$. Thus, $(S + I, +, \cdot)$ is a subring of $(R, +, \cdot)$.
        \end{proof}

        \item[(b)]
        \begin{theorem}\label{thm:ideal_subset}
            Prove: $I \trianglelefteq S + I$.
        \end{theorem}
        \begin{proof}
            Consider $a \in I$ and $b' \in S + I$ such that $b' = s + b$ where $s \in S$, $b \in I$. Observe
            \[
            ab = as + ab, \quad as \in I, \quad ab \in I \implies ab \in I
            \]
            and $b'a \in I$ similarly.
        \end{proof}

        \item[(c)]
        \begin{theorem}\label{thm:intersection_subring}
            Show: $S \cap I \trianglelefteq S$.
        \end{theorem}
        \begin{proof}
            As $(I, +) \trianglelefteq R$,
            \[
            (S \cap I, +) \trianglelefteq S \quad \text{and} \quad 0 \in S \cap I \implies S \cap I \neq \emptyset.
            \]
            If $a \in S \cap I$ and $s \in S$, then $a \in I$ and $s \in S \implies as \in I$. Now, $a \in S$ and $s \in S \implies as \in S$.
        \end{proof}
    \end{enumerate}
\end{enumerate}

% --- Page 10 ---
---

\[
\text{Hence, as } s \in S \cap I \implies \varphi(s) = 0 + I \text{ as } s \in S \cap I.
\]

\[
\implies S \cap I \subseteq \ker \varphi.
\]

(cd) Prove:
\[
\frac{S + I}{I} \cong \frac{S}{S \cap I}
\]

**Proof.**

1. Define $\varphi: S \to \frac{S + I}{I}$ by $\varphi(s) = s + I$.

2. **Well-definedness:**
   If $s_1 = s_2$, then $s_1 + I = s_2 + I$.

3. $\varphi$ is a ring homomorphism:
   \[
   \varphi(s + s') = s + s' + I = (s + I) + (s' + I) = \varphi(s) + \varphi(s'),
   \]
   \[
   \varphi(ss') = ss' + I = (s + I)(s' + I) = \varphi(s)\varphi(s').
   \]

4. $\varphi$ is surjective (by construction).

5. $\ker \varphi = S \cap I$.

   If $s \in S \cap I$, then $\varphi(s) = 0 + I$ as $s \in S \cap I$. Hence, $S \cap I \subseteq \ker \varphi$.

   If $s \in \ker \varphi$, then $\varphi(s) = I$, now $s \in S$ and $s \in I \implies s \in S \cap I$.

   \[
   \therefore \quad S \cap I = \ker \varphi.
   \]

By the First Isomorphism Theorem,
\[
\frac{S + I}{I} \cong \frac{S}{S \cap I}.
\]

---

% --- Page 11 ---
Given: $D \in \mathbb{Q}$ be a rational number and is NOT a perfect square.
\[ R := \{ a + b\sqrt{D} \mid a, b \in \mathbb{Q} \} \]

(a)
\begin{proof}
Observe, $R \neq \emptyset$ as $0 \in R$.

Consider $x, y \in R$, where $x = a_1 + b_1 \sqrt{D}$ and $y = a_2 + b_2 \sqrt{D}$; $a_1, b_1, a_2, b_2 \in \mathbb{Q}$.

\[ x - y = (a_1 - a_2) + (b_1 - b_2) \sqrt{D} \in R \text{ as } a_1 - a_2 \in \mathbb{Q} \text{ and } b_1 - b_2 \in \mathbb{Q}. \]

Thus, $(R, +)$ is a subgroup of $(\mathbb{C}, +)$.

Now, look at
\[ xy = (a_1 + b_1 \sqrt{D})(a_2 + b_2 \sqrt{D}) = a_1 a_2 + b_1 b_2 D + (a_2 b_1 + a_1 b_2) \sqrt{D}, \]
and $a_1 a_2 + b_1 b_2 D, a_2 b_1 + a_1 b_2 \in \mathbb{Q}$.

Hence, $xy \in R$ and $(R, +, \cdot)$ is a subring of $(\mathbb{C}, +, \cdot)$.
\end{proof}

(b)
\begin{theorem}\label{thm:nonzero}
Show: If $(a + b\sqrt{D}) \neq 0$, then show that $a^2 - b^2 D \neq 0$.
\end{theorem}

\begin{proof}
If $(a + b\sqrt{D}) \neq 0$, then $a \neq 0$ and $b \neq 0$.

Hence, $a - b\sqrt{D} \neq 0$. Now, $(a + b\sqrt{D})(a - b\sqrt{D}) \neq 0$

\[
\Rightarrow a^2 - b^2 D + 0 \cdot (\sqrt{D}) \neq 0.
\]
\end{proof}

(c)
\begin{theorem}\label{thm:field}
Prove: $R$ is a field.
\end{theorem}

\begin{proof}
Observe $R$ is a commutative, unital ring.

Moreover, $R$ is a division ring as for $a + b\sqrt{D} \neq 0$, then
\[
\frac{a - b\sqrt{D}}{a^2 - b^2 D}
\]
is the inverse with respect to multiplication.

Now, $R$ is a commutative, division ring, which is a field.
\end{proof}

% --- Page 12 ---
Given the content in the image, 
Given $R$ is a unital ring,

\[
Z(CR) := \{ a \in R \mid ab = ba \text{ for all } b \in R \}.
\]

\begin{proof}
Observe that $Z(CR) \neq \emptyset$ as $1 \in Z(CR)$.

Now, consider $a, a' \in Z(CR)$. For any $b \in R$,
\[
(a - a')b = ab - a'b = ba - ba' = b(a - a') \implies (Z(CR), +) \text{ is a subgroup of } R.
\]

For any $b \in R$, observe $aa'b = aba' = baa'$ $\implies aa' \in Z(CR)$.

Hence, $Z(CR)$ is a subring.
\end{proof}

\end{document}